
\vspace{-0.2cm}

\section{Conclusions}



Most prior work improves image generators under a fixed dataset regime. We take a different path: we scale the language itself by training on long structured captions that encode far richer visual and compositional detail than any prior open-source approach. This shift changes learning dynamics and behavior: long captions accelerate convergence and yield native disentanglement, so single-attribute edits (lighting, depth of field, expression) modify only the intended factor while the rest remains stable. To make this feasible, we introduced DimFusion, a novel architecture that integrates intermediate and final LLM representations without increasing token length. And because thousand-word prompts outstrip human evaluators, we proposed Text-as-a-Bottleneck Reconstruction (TaBR)—an image-grounded protocol to measure expressiveness and controllability. All components culminate in \modelname{}, a large-scale text-to-image model that outperforms larger baselines, demonstrating that caption scale and structure are levers for controllable generation. Looking ahead, our results suggest a broader design space: we can build intermediate languages that unlock new user interfaces. Users continue to write natural language, which is translated into a structured intermediate dialect that enables new forms of creative control. This paves the way for predictable graphics primitives atop image generation models.





\section{Ethical Statement}
\label{sec:ethics}



Most text-to-image models are trained on data scraped online without permission or licenses from the rightful owners, a practice that violates intellectual property rights and yields models that cannot guarantee the copyright status of their outputs. In contrast, our model is trained solely on licensed data, substantially reducing the risk of infringement. We believe the AI ethics community should prioritize licensed, consent-based training, as responsible data use is essential for a sustainable AI ecosystem. Complementing this stance, our work investigates how structured synthetic captions and LLM fusion can improve fidelity and controllability: by strengthening prompt adherence and semantic grounding, our approach can help reduce unintended biases and misalignment. At the same time, greater controllability introduces dual-use risks, as it may also be exploited to generate misleading, or inappropriate content more reliably. We therefore advocate openly studying both the benefits and risks of these techniques to develop effective safeguards.